\begin{vidu}%[2P5-5.2B][Loai1]
Một ống dây có hệ số tự cảm là 0,01 H. Khi có dòng điện chạy qua ống dây có năng lượng 0,08 J. Cường độ dòng điện chạy qua ống dây bằng
\choice
{1 A}
{2 A}
{3 A}
{\True 4 A}
\loigiai{ Năng lượng mà nguồn điện cung cấp cho ống dây:
$W=\dfrac{1}{2}LI^2\Rightarrow I=\sqrt{{\dfrac{2W}{L}}}=4\ A. $}
\end{vidu}

\begin{vidu}%[2P5-5.2K][Loai1]
Một ống dây dài 40 cm có tất cả 800 vòng dây. Diện tích tiết diện ống dây là $10\ cm^2$. Cường độ dòng điện qua ống tăng từ 0 đến 4 A. Hỏi nguồn điện đã cung cấp cho ống dây một năng lượng bằng
\choice
{\True $1,6.10^{-2}$ J}
{$1,8.10^{-2}$ J}
{$2.10^{-2}$ J}
{$2,2.10^{-2}$ J}
\loigiai{\begin{itemize}
\item Độ tự cảm của ống dây: $L=4\pi.10^{-7}\dfrac{N^2S}{\ell}=4\pi.10^{-7}\dfrac{{800}^2{.10.10}^{-4}}{0,4}\approx 2,01.10^{-3}\ H $.
\item Năng lượng mà nguồn điện cung cấp cho ống dây:
\begin{align*}
W=\dfrac{1}{2}L\left|{I_2^2-I_1^2}\right|=\dfrac{1}{2}2,01.10^{-3}\left|{16-0}\right|=1,6.10^{-2}\ J.
\end{align*}
\end{itemize}
}
\end{vidu}
\begin{vidu}%[2P5-5.2K][Loai1]
	Dòng điện qua một ống dây không có lõi sắt biến đổi đều theo thời gian, trong 0,01 s cường độ dòng điện tăng đều từ 1 A đến 2 A thì suất điện động tự cảm trong ống dây là 20 V. Tính hệ số tự cảm của ống dây và độ biến thiên năng lượng của từ trường trong ống dây.
	\choice
	{0,2 H; 0,5 J}
	{0,1 H; 0,2 J}
	{0,3 H; 0,4 J}
	{\True 0,2 H; 0,3 J}
	\loigiai{
		\begin{itemize}
			\item Hệ số tự cảm của ống dây là: $L = \dfrac{{e_{tc}.\Delta t}}{\Delta I} = \dfrac{{20.\left( {2 - 1} \right)}}{0,01} = 0,2$ H.
			\item Độ biến thiên năng lượng của từ trường trong ống dây là: $W= \dfrac{1}{2}.0,2.\left( {{2^2} - {1^2}} \right) = 0,3$ J.
		\end{itemize}
	}
\end{vidu}